\subsection{An actual quadratic character detecting the CM discrepancy}

\begin{statement}{E.65}{Proved}{The CM parity is a continuous character of the full gentle group}
Put $\mathcal G=\widehat{\mathrm{GT}}_0$ and
$\mathcal H=\ker(\lambda:\mathcal G\to\widehat{\mathbf Z}^{\times})$.
Retain the positive generators and actual pure frames of E.58 and E.61.
For $F=(1,f)\in\mathcal H$, write
\[
 \Phi_F(s)=\beta_Fs\beta_F^{-1},\qquad
 \Phi_F(r)=A_FrA_F^{-1},\qquad
 A_F,\beta_F\in\widehat F,
 \qquad (\beta_F^{-1})^{\rm ab}=(\rho_F,\rho_F).
 \tag{E.65.1}
\]
Both frames are unique. Use the signed map $q_{12}$ of E.61.3 and set
\[
 \begin{gathered}
 U=\begin{pmatrix}1&1\\0&1\end{pmatrix},\quad
 L=\begin{pmatrix}1&0\\-1&1\end{pmatrix},\quad
 S=ULU,\quad R_+=UL,\quad R_0=-R_+,\\
 V_F=U^{8\rho_F},\qquad
 T_F=q_{12}(A_F)^{-1}V_F,\qquad
 W_F=q_{12}(\beta_F)^{-1}V_F
                    \quad\text{over }\mathbf Z/12.
 \end{gathered}
 \tag{E.65.2}
\]
Here $R_0^3=I$, $R_+^3=-I$, and
$\ker(R_0-I)=\ker(R_++I)=\mathbf F_3(1,-1)$ at three.
The matrices $T_F,W_F$ are both $I$ modulo four. Their reductions
at three give actual continuous homomorphisms
\[
 \begin{split}
 T:\mathcal H&\longrightarrow
 C_{\mathrm{SL}_2(\mathbf F_3)}(R_0)
       =\langle -I\rangle\times\langle R_0\rangle\simeq C_6,\\
 W:\mathcal H&\longrightarrow
 C_{\mathrm{SL}_2(\mathbf F_3)}(S)
       =\langle S\rangle\simeq C_4.
 \end{split}
 \tag{E.65.3}
\]
Define $\epsilon_3(T)$ by
$T(1,-1)=(-1)^{\epsilon_3(T)}(1,-1)$, and define
$\epsilon_2(S^j)=j\bmod2$. Then
\[
       \delta(F)=\epsilon_3(T_F)+\epsilon_2(W_F)
                    \quad\text{in }\mathbf F_2
 \tag{E.65.4}
\]
is a continuous character. It is exactly the extra-CM parity of E.1,
including its global-sign convention.

This character has a unique extension
\[
 \Delta:\mathcal G\longrightarrow\mathbf F_2,
 \qquad \Delta|_{G_{\mathbf Q}}=0,
 \qquad \Delta|_{\mathcal H}=\delta.
 \tag{E.65.5}
\]
For an arbitrary $F\in\mathcal G$, the following is an exact,
choice-independent formula: choose an actual Galois element
$\sigma$ with $\chi(\sigma)=\lambda(F)$ and put
\[
                    \Delta(F)=\delta(F\sigma^{-1}).
 \tag{E.65.6}
\]
The full gentle element constructed in E.58 has $\Delta=1$.
Thus $\ker\Delta$ is an actual closed subgroup of index two in
$\mathcal G$, containing the Galois image.

The remaining Task E.3 is precisely the computation of the image of
the displayed restriction
\[
              \Delta|_{\widehat{\mathrm{GT}}}:
                         \widehat{\mathrm{GT}}\longrightarrow\mathbf F_2.
 \tag{E.65.7}
\]
The construction of this finite character does not assign that image.
\end{statement}
\begin{proof}
We establish the finite covariance, the sign maps, and their transport
under arbitrary multipliers before extending the character.

\paragraph{The actual finite covariance.}
The action of $\mathcal G$ on $\widehat B$ is the quotient of its
standard action on $\widehat B_3$:
\[
 u\longmapsto u^\lambda,\qquad
 v\longmapsto f(u^2,v^2)^{-1}v^\lambda f(u^2,v^2).
 \tag{E.65.8}
\]
The two gentle hexagons give the braid relation and the image
$(uv)^3\mapsto(uv)^{3\lambda}$; the inverse gentle pair supplies
the inverse map. Thus this is the actual automorphism and central
quotient used throughout E.58--E.62. On $B$ one has
$s=uvu$, $r=uv$, and $sr=u^{-1}$, where the equalities are taken
modulo the central full twist. In particular the multiplier-one
action fixes $sr$ literally.

The torsion centralizers from E.58 give
$C_{\widehat B}(s)\cap\widehat F=
C_{\widehat B}(r)\cap\widehat F=1$.
They prove uniqueness in \textup{(E.65.1)}. They also prove
continuity: conjugation from the compact pure group to the indicated
conjugacy class is a continuous bijection onto a Hausdorff space.

E.61.4 gives the signed identity
\[
 q_{12}(f)=
 \begin{pmatrix}1&-8\rho_F\\-8\rho_F&1+64\rho_F^2\end{pmatrix}.
 \tag{E.65.9}
\]
Direct multiplication, in the specified modulus twelve, gives
\[
 V_FS V_F^{-1}=S q_{12}(f),\qquad
 V_FR_0V_F^{-1}=q_{12}(f)^{-1}R_0.
 \tag{E.65.10}
\]
These are exactly the images of $\Phi_F(s)=sf$ and
$\Phi_F(r)=f^{-1}r$. Hence
$q\Phi_F=\operatorname{Ad}(V_F)q$ on the actual quotient
$q:\widehat B\to\mathrm{SL}_2(\mathbf Z/12)/\{\pm I\}$.
Its restriction to $\widehat F$ is also the signed equality
$q_{12}\Phi_F=\operatorname{Ad}(V_F)q_{12}$: both maps have
the same modulo-four value on $\widehat F^{\rm ab}$, and
$-I$ does not belong to the displayed $V_4$ in E.61.9.
This argument applies to every full gentle pair and uses no
four-strand action.

\paragraph{The integral sign that fixes the lift.}
Use the free class-two quotient with coordinates
$x^a y^b[x,y]^c$ over $\widehat{\mathbf Z}$. The multiplier-one
three-cycle equation sends the central image $[x,y]^c$ of $f$
to $[x,y]^{3c}=1$. Since $\widehat{\mathbf Z}$ has no
three-torsion, $c=0$.

Put $h=\beta_F^{-1}$ and write its class-two image as
$x^{\rho_F}y^{\rho_F}[x,y]^d$. The actual interchange
$\theta(x)=y$, $\theta(y)=x$ gives
\[
 \theta(h)^{-1}h=[x,y]^{\rho_F^2+2d}.
 \tag{E.65.11}
\]
This is $f$ by E.61.2. Therefore $\rho_F^2+2d=0$, so
$\rho_F\in2\widehat{\mathbf Z}$. The image of $\widehat F$
modulo four is E.61.9's elementary four-group. Consequently
$q_4(\beta_F)=I$.

Also $V_F=I$ modulo four. The first pure-frame relation for $r$
and \textup{(E.65.10)} imply that $q_4(A_F)$ centralizes
$R_0$ in that four-group. This centralizer is trivial: substitution
of its four displayed matrices gives only $I$. Thus $q_4(A_F)=I$.
This proves the two assertions modulo four without changing any
global sign.

The projective frame equalities now lift with this same sign to
give $T_FR_0=R_0T_F$ and $W_FS=SW_F$ modulo twelve. One may
also distinguish the sign for $R_0$ by its trace $-1$, and the
sign for $S$ by its reduction modulo four. Solving the two
centralizer equations at three gives the exact groups in
\textup{(E.65.3)}.

\paragraph{The frame multiplication law.}
Use $FG$ for the composition $\Phi_F\Phi_G$.
Uniqueness of the pure frames gives
\[
 A_{FG}=\Phi_F(A_G)A_F,\qquad
 \beta_{FG}=\Phi_F(\beta_G)\beta_F.
 \tag{E.65.12}
\]
The multiplier-one action is the identity on $\widehat F^{\rm ab}$,
so $\rho_{FG}=\rho_F+\rho_G$ and $V_{FG}=V_FV_G$.
Substituting the signed covariance in \textup{(E.65.12)} gives
$T_{FG}=T_FT_G$ and $W_{FG}=W_FW_G$. This proves the asserted
homomorphisms. Both characters in \textup{(E.65.4)} are literal
characters of their indicated finite cyclic groups.

The sign comparison is the one in E.1: a matrix $T_F$ that is $I$
modulo four must be multiplied by its sign on the rational
three-torsion line to choose the global sign fixing that point.
Its resulting four-torsion action changes $r_0=(3/4)^{1/4}$
by this sign. On the Gaussian side the class of $S^j$ modulo
$\{\pm I\}$ changes $r_0$ by $(-1)^j$. These are the two
actual restrictions to
\[
 \mathbf Q(E_3[4])=\mathbf Q(i,r_0)
              =\mathbf Q\bigl(x(E_4[3])\bigr)
 \tag{E.65.13}
\]
proved by the explicit division-point formulas in LM7 D.25 and
retained in E.1. Thus \textup{(E.65.4)} is exactly their difference.

\paragraph{Transport under nontrivial multipliers.}
Here we retain the full source maps, including the non-pure frame
at the order-three factor. For $G\in\mathcal G$ let
$\varepsilon_G\in\{1,-1\}$ be its multiplier modulo three and
$b_G=(1-\varepsilon_G)/2$. The normalized frames are
$\beta_G\in\widehat F$ and $A_G\in\widehat F s^{b_G}$, with
\[
 \Phi_G(s)=\beta_Gs\beta_G^{-1},\qquad
 \Phi_G(r)=A_Gr^{\varepsilon_G}A_G^{-1}.
 \tag{E.65.14}
\]
The action on $S_3$ is the identity, since the two braid generators
have odd multiplier and $f$ is pure. Thus the indicated cosets
are exactly the required conjugator cosets. Multiplication by a
right power of the respective torsion generator chooses that
coset, and the torsion-centralizer theorem proves uniqueness.
It follows that \textup{(E.65.12)} continues to hold for these
frames: its right side has the prescribed $S_3$ coset and the
correct torsion power.

For clarity, the needed congruence intertwiner can be obtained
without identifying different Kummer normalizations. On
$\widehat B_3$ use the signed representation $u\mapsto U$,
$v\mapsto L$ over each $\mathbf Z_\ell$. Its composition with
\textup{(E.65.8)} sends $u$ to $U^\lambda$ and $v$ to a
trace-two matrix satisfying the braid relation with $U^\lambda$.
Direct solution of that relation over the domain $\mathbf Z_\ell$
gives either the abelian pair or a unique intertwiner
\[
 \mathcal V_G=\begin{pmatrix}\lambda(G)&t_G\\0&1\end{pmatrix},
 \qquad
 q_\ell\Phi_G=\operatorname{Ad}(\mathcal V_G)q_\ell.
 \tag{E.65.15}
\]
The abelian pair is excluded because precomposition with an
automorphism preserves the nonabelian image of $\langle U,L\rangle$.
Explicitly the second matrix is
$\left(\begin{smallmatrix}1-u&\lambda u^2\\-1/\lambda&1+u\end{smallmatrix}\right)$,
and $t_G=\lambda u$. This is the calculation of
\cite[Lemmas 4.1--4.3]{DE_Marin2008}, with the displayed normalization.
The centralizer of both $U$ and $L$ consists of scalars; fixing the
lower-right entry to one makes the intertwiner unique. Thus
$\mathcal V_{GH}=\mathcal V_G\mathcal V_H$.

Taking the product over $\ell$ gives an actual matrix over
$\widehat{\mathbf Z}$. For $F\in\mathcal H$ its reduction modulo
twelve is exactly $V_F$: \textup{(E.65.10)} supplies that covariance,
and two normalized unipotent shears with the same covariance
commute with both $U$ and $L$, hence are equal.

Use projective matrices for the non-pure $A_G$. The actual
quantities
$\mathcal T_G=q(A_G)^{-1}\mathcal V_G$ and
$\mathcal W_G=q(\beta_G)^{-1}\mathcal V_G$ are projective
representations by the same calculation as
\textup{(E.65.12)}. They normalize respectively $\langle R_0\rangle$
and $\langle S\rangle$. Hence for $F\in\mathcal H$,
\[
 T_{GFG^{-1}}=\mathcal T_GT_F\mathcal T_G^{-1},\qquad
 W_{GFG^{-1}}=\mathcal W_GW_F\mathcal W_G^{-1}.
 \tag{E.65.16}
\]
These equalities are signed modulo twelve: both sides are $I$
modulo four, which chooses the same global lift.

Every normalizer of $\langle R_0\rangle$ at three preserves
$\mathbf F_3(1,-1)$ and therefore preserves its scalar character
under conjugation. Every normalizer of $\langle S\rangle$ sends
$S$ to $S$ or $S^{-1}$ and therefore preserves its exponent
modulo two. Consequently
\[
                         \delta(GFG^{-1})=\delta(F).
 \tag{E.65.17}
\]
This proves the required invariance for all multipliers using
actual conjugations.

\paragraph{Galois vanishing and the full extension.}
For a Galois element of multiplier one, both restrictions in
\textup{(E.65.13)} are its action on the same field, so their
parities agree. One can use only the already fixed E.1 matrix
normalization in this step: on the kernel of $\rho\bmod3$ one
has $V_F=I$ modulo twelve, and E.1 gives $\delta=0$ there.
The quotient of the Galois multiplier-one group by that kernel
has order dividing three, so a character to $\mathbf F_2$ must
vanish on the entire group. This also proves the claim without
changing the shear convention in any arithmetic formula.

The cyclotomic character of $G_{\mathbf Q}$ is onto
$\widehat{\mathbf Z}^{\times}$. Hence every $F\in\mathcal G$
has a decomposition $F=H\sigma$ with $H\in\mathcal H$ and
$\sigma\in G_{\mathbf Q}$. If two choices have the same multiplier,
their ratio belongs to $G_{\mathbf Q}\cap\mathcal H$, on which
$\delta=0$. Thus \textup{(E.65.6)} is independent of the choice.
The conjugation identity \textup{(E.65.17)} proves directly that
it is a homomorphism. Continuity follows from the actual compact
surjection
\[
       \mathcal H\times G_{\mathbf Q}\longrightarrow\mathcal G,
                           \qquad(H,\sigma)\longmapsto H\sigma:
 \tag{E.65.18}
\]
the continuous function $\delta(H)$ is constant on every fibre,
and a continuous map from a compact space onto a Hausdorff space
is a quotient map. This also proves uniqueness in
\textup{(E.65.5)}.

E.58 has $T=-I$ at three in the lift equal to $I$ at four, and
$W=I$, so its value is $1$. Finally multiplication by an actual
Galois element of the same multiplier preserves $\Delta$ and
reduces any full-GT element to $H\in\mathcal H$. Take the actual
group element $H^3$. The proved character and frame laws give
$\Delta(H^3)=\Delta(H)$ and $\rho_{H^3}=3\rho_H$.
Consequently $V_{H^3}=I$ and $q_{12}(f_{H^3})=I$ modulo twelve.
Its action fixes $Q$ and therefore preserves $N_B$ and its
characteristic subgroup $N_B'N_B^6$. It induces an actual
automorphism of the specified $G_6$, with the E.1 pure frames
and the same parity. Thus every nonzero value of the restricted
character yields a point in the actual extra E.1 CM matrix
fibre. The reverse implication is \textup{(E.65.4)}.
This proves the exact reduction to \textup{(E.65.7)} without
an additional Kummer-surjectivity assumption.

\paragraph{Finite certificate and precise source scope.}
The independent standard-library program
\texttt{a\_cm\_parity\_character.py} verifies all twelve finite
covariances, both complete Cartans, all $24^2$ products, and
every conjugation by their full normalizers in
$\mathrm{GL}_2(\mathbf F_3)$. It also checks the integral
class-two coordinate identity and the central-difference character
on the actual dihedral splitting-field law for $X^4-3/4$.
All outputs are retained. No finite test in that program assumes
or proves the image in \textup{(E.65.7)}.
\end{proof}
